\documentclass[UTF8,11pt,openany]{ctexrep}

\usepackage[a4paper,margin=2.4cm]{geometry}
\usepackage{amsmath,amssymb,mathtools}
\usepackage{booktabs}
\usepackage{enumitem}
\usepackage{array}
\usepackage{xcolor}
\usepackage[colorlinks=true,linkcolor=blue,urlcolor=blue,citecolor=blue]{hyperref}

\title{《花砖物语：琉璃之光》的有限时域随机策略模型}
\author{}
\date{模型设定与策略搜索记录}

\newcommand{\ind}{\mathbb{I}}
\newcommand{\xvec}{\boldsymbol{x}}
\newcommand{\mvec}{\boldsymbol{m}}

\begin{document}
\maketitle

\begin{abstract}
本文建立一个受《花砖物语：琉璃之光》（\textit{Azul: Stained Glass of
Sintra}）启发的有限时域随机决策模型。模型保留三个核心策略因素：玻璃匠只能
向右移动、回到左侧需要消耗一个完整行动，以及右侧已经完成的窗户会被左侧窗户
反复计分。场面花色带来的选择限制被抽象为每次决策前可观察的随机行动Mask。

第一章固定模型设定、状态、奖励与Bellman方程，作为后续实验的共同背景。第二章
按实际发生的顺序记录策略搜索：从精确动态规划的可扩展性问题，到可解释参数化
策略，再到由讨论中的反例触发的特征修正。文中刻意保留被后续结果推翻的中间
结论，并同时说明每轮方法的局限，避免把迭代过程写成事后看来必然正确的路线。
当前纯最高分策略在 $H=45$ 的独立样本验证中平均得到约 $45.57$ 分；采用更平衡
的重置阈值时平均仍有 $45.57$ 分、完成 $10.49$ 面且总重置降至 $10.47$ 次。
这些策略仍是近似策略而非全局最优解。
\end{abstract}

\tableofcontents

\chapter{初始简化设定与问题背景}

\section{问题设定}

场上有八扇从左到右排列的窗户，编号为 $i\in\{1,\ldots,8\}$。每扇窗户可以
完成两次，对应原游戏中彩窗条的正反两面；每一面需要三次成功填充。因此，一扇
窗户总共需要六次成功填充才会完全移出游戏。

游戏使用固定行动时域
\[
H=45.
\]
填充、主动重置和被迫空过都会消耗一次行动。当四十五次行动用完，或者八扇窗户
全部完成两次时，游戏结束。

每完成一面，基础得分为一分。如果完成的是第 $i$ 扇窗户，则右侧每一扇至少完成
过一次的窗户额外提供一分。同一扇右侧窗户即使完成了两次，在右侧加分中仍然只
计算一次。

\section{状态表示}

定义
\[
x_i\in\{0,1,\ldots,6\}
\]
为第 $i$ 扇窗户累计完成的填充次数。该窗户已经完成的面数为
\[
c_i=\left\lfloor\frac{x_i}{3}\right\rfloor\in\{0,1,2\},
\]
当前一面的填充进度为
\[
q_i=x_i\bmod 3.
\]
当 $x_i=6$ 时，该窗户已经完成两次，不可继续填充。

令 $g\in\{1,\ldots,8\}$ 表示玻璃匠当前所在的物理位置，$h$ 表示剩余行动
次数。随机Mask产生前的完整状态为
\[
s=(\xvec,g,h),\qquad \xvec=(x_1,\ldots,x_8).
\]
由于得分是可加的，而且历史累计分数不会影响未来状态转移，因此累计得分不需要
进入状态。在任意非终止状态，定义最左侧仍然存在的窗户为
\[
\ell(\xvec)=\min\{i:x_i<6\}.
\]

\section{行动前可观察的随机Mask}

每次行动开始时，先为每扇尚未完全完成的窗户生成二元Mask：
\[
m_i\sim\operatorname{Bernoulli}(p(q_i)),
\qquad
p(q_i)=
\begin{cases}
0.8, & q_i=0,\\
0.6, & q_i=1,\\
0.4, & q_i=2.
\end{cases}
\]
也就是说，当前一面还差三次、两次和一次填充时，本行动可填充的概率分别为
$80\%$、$60\%$ 和 $40\%$。第一面完成后，第二面的进度重新归零，因此概率也
重新回到 $80\%$。完全完成的窗户不再生成Mask。

基准模型假设各窗户的Mask在给定状态后条件独立。玩家在选择本行动之前观察完整
的 $\mvec$。这一时间顺序是模型的关键：玩家不是先选择窗户再判断成功，而是先
看到哪些窗户可选，再从合法行动中选择期望价值最高者。

\section{行动与状态转移}

\subsection{填充行动}

观察到 $\mvec$ 后，填充第 $i$ 扇窗户是合法的，当且仅当
\[
i\ge g,\qquad x_i<6,\qquad m_i=1.
\]
即玩家只能填充玻璃匠当前所在位置或其右侧的可用窗户。执行填充后，
\[
x_i'=x_i+1,\qquad g'=i,
\]
其余窗户的填充状态保持不变。

\subsection{重置行动}

重置消耗一个完整行动，不产生即时得分，并令
\[
g'=\ell(\xvec).
\]
下一次行动会重新生成一组Mask。当玻璃匠不在最左侧有效窗户时，无论当前是否有
合法填充行动，玩家都可以主动重置。

如果玻璃匠当前位置及其右侧没有任何合法填充行动，则必须重置。如果此时玻璃匠
已经位于最左侧有效窗户，重置等价于被迫空过：状态不变，但仍然消耗一个行动。
当玻璃匠已经在最左侧且至少存在一个合法填充行动时，不允许主动空过。

\section{即时奖励}

普通填充的即时奖励为零。只有当填充完成一面时才产生奖励，即行动前满足
\[
x_i\in\{2,5\}.
\]
此时即时得分为
\[
r_i(\xvec)
=1+\sum_{j=i+1}^{8}\ind\{x_j\ge3\}.
\]
指标函数 $\ind\{x_j\ge3\}$ 保证右侧窗户第一次完成后开始提供一分加成，第二次
完成不会将加成提高到两分。当前窗户的第一次和第二次完成都会触发上述完整计分。

例如，若第八扇窗户已经至少完成一次，则第七扇窗户每次完成都得到两分：一分
基础分加上第八扇窗户提供的一分。即使第八扇窗户已经完成两次，第七扇窗户仍然
只得到两分。

若不存在行动时域限制，并且能够按照完美的从右到左顺序将所有窗户完成两次，
理论最高得分为
\[
2\sum_{k=1}^{8}k=72.
\]
但在四十五次行动内不可能完成全部十六面，因为仅成功填充就需要四十八次行动，
尚未计算重置与被迫空过的成本。

\section{终局价值}

当 $h=0$ 或所有窗户均满足 $x_i=6$ 时，过程终止。未完成一面的残留进度没有
任何终局价值。因此，终局时处于 $x_i\in\{1,2,4,5\}$ 的窗户不会获得额外分数；
总分完全由游戏过程中已经触发的完成奖励组成。

在四十五次行动下，完成面数的绝对上界为
\[
\left\lfloor\frac{45}{3}\right\rfloor=15.
\]
只有在没有任何重置、空过或终局残留进度时才能达到该上界。因此，平均完成次数
约为九至十次是本模型希望覆盖的实际区间。

\section{有限时域Bellman方程}

令 $V_h(\xvec,g)$ 表示剩余 $h$ 次行动、尚未生成当前Mask时的最优期望未来
得分。观察到Mask后，令 $\mathcal{A}(\xvec,g,\mvec)$ 为所有合法填充和重置
行动的集合，则
\[
V_h(\xvec,g)
=\mathbb{E}_{\mvec\mid\xvec}
\left[
\max_{a\in\mathcal{A}(\xvec,g,\mvec)}
\left\{
r(\xvec,a)+V_{h-1}(\xvec',g')
\right\}
\right],
\]
边界条件为
\[
V_0(\xvec,g)=0.
\]
期望位于最大化算子之外，表示玩家先观察随机Mask，再选择当前最优行动。

由于各窗户Mask条件独立，在计算一个固定状态的期望最大值时，不必显式枚举
$2^8$ 种Mask。可以先计算所有填充行动的续接价值，将其从高到低排序，再根据
“所有更优行动均不可用、当前行动可用”的概率逐项累加，从而精确边缘化随机
Mask。这一化简消除了单个状态内部的Mask枚举，但没有消除跨行动状态数量的指数
增长；该区别会在第二章成为选择搜索方法的关键。

\section{希望提取的策略结论}

最优策略是映射
\[
\pi^*(\xvec,g,h,\mvec)
\longrightarrow
\{\text{填充窗户 }i,\text{重置}\}.
\]
本研究不只希望得到一个难以解释的策略表，还希望提取以下可复用的结论：
\begin{enumerate}[leftmargin=*]
    \item 初始状态下，如果多个窗户同时可用，第一步优先填充哪一扇；
    \item 随机Mask下，真实的开局选择分布与条件优先级；
    \item 何时值得牺牲一个完整行动重置，而不是继续向右；
    \item “尽早激活右侧加分”与“保留左侧行动灵活性”之间的权衡；
    \item 第一面负责建立加分结构、第二面负责兑现分数时，策略应在何时切换阶段。
\end{enumerate}

\chapter{策略搜索过程与讨论驱动的迭代}

\section{精确动态规划的第一次尝试}

Bellman方程给出了精确定义，但直接求解并不轻量。忽略可达性约束，仅Mask生成前
的状态空间上界就达到
\[
7^8\times 8\times 46=2{,}121{,}446{,}768.
\]
独立Mask可以用排序后的动作价值精确积分，因此 $2^8$ 种Mask并不是主要瓶颈；
真正困难的是八扇窗户进度、玻璃匠位置和剩余行动数共同造成的状态爆炸。

一次朴素的稀疏记忆化动态规划在较短的 $H=36$ 测试中，已经访问超过
$4.2\times10^8$ 个状态并产生超过 $2.1\times10^9$ 次递归调用，仍无法在合理
资源内完成。这个结果不能证明所有精确算法都不可行，却足以说明当前实现不适合
直接扩展到 $H=45$。因此搜索路线转向“可解释参数化策略加大规模蒙特卡洛评估”，
并保留未来使用状态聚合、界限剪枝或近似动态规划的空间。

\section{第一版可解释参数化策略}

第一版策略对每个当前合法的填充动作计算线性效用
\[
u_i(s,h)=\boldsymbol{w}^{\mathsf T}\boldsymbol{\phi}(s,i,h),
\]
然后选择效用最高的窗户。为完整写出初始的十四个特征，先定义以下辅助量：
\[
z_i=\frac{i-1}{7},\qquad
d_i=\ind\{q_i=2\},\qquad
B_i(\xvec)=1+\sum_{j=i+1}^{8}\ind\{x_j\ge3\},
\]
其中 $z_i\in[0,1]$ 是从左到右的位置（程序内部使用从零开始的编号），$d_i$
表示本次填充是否完成当前一面，$B_i$ 是若本次完成时能够取得的未归一化即时
分数。再定义
\[
C_i(\xvec)=\sum_{k=1}^{i-1}\max\{0,2-c_k\},\qquad
F_i(\xvec)=\frac{1}{8}\sum_{j=i}^{8}\ind\{x_j<6\}.
\]
$C_i$ 是第 $i$ 扇窗户左侧尚未兑现的完成面总数，用来近似第一次完成第 $i$ 扇
后能够激活的未来加分机会；$F_i$ 是从位置 $i$ 开始的有效窗户比例，用来近似
填充后仍然保留的右侧行动灵活性。

按实际实现中的顺序，特征向量为
\[
\begin{aligned}
\boldsymbol{\phi}(s,i,h)=\bigg(&
z_i,\ \frac{q_i}{2},\ \ind\{c_i=0\},\ d_i,\
\frac{d_iB_i(\xvec)}{8},\
\frac{d_i\ind\{c_i=0\}C_i(\xvec)}{14},\
\ind\{i=g\},\ F_i(\xvec),\\
&d_i\frac{h}{H},\
\frac{q_i}{2}\left(1-\frac{h}{H}\right),\
\ind\{c_i=1\},\ \ind\{x_i<6\},\
\ind\{i=8\},\ z_i^2
\bigg)^{\mathsf T}.
\end{aligned}
\]
逐项含义如下：
\begin{enumerate}[label={$\phi_{\arabic*}$：},leftmargin=*]
    \item $z_i$，线性位置特征；数值越大表示窗户越靠右。
    \item $q_i/2$，当前一面的归一化填充进度。
    \item $\ind\{c_i=0\}$，当前是否仍处于第一面。
    \item $d_i=\ind\{q_i=2\}$，本次填充是否立即完成一面。
    \item $d_iB_i(\xvec)/8$，本次完成能够得到的归一化即时分数；没有立即
    完成时该特征为零。
    \item $d_i\ind\{c_i=0\}C_i(\xvec)/14$，第一次完成当前窗户时对左侧
    未来完成机会的激活价值。分母 $14=2\times7$ 是该计数的理论最大值。
    \item $\ind\{i=g\}$，是否停留在玻璃匠当前所在窗户。
    \item $F_i(\xvec)$，当前位置及其右侧仍有效的窗户比例，即可达性灵活度代理。
    \item $d_i h/H$，“立即完成”与剩余时域的交互；同样的完成在较早阶段
    允许更多后续结构收益。
    \item $(q_i/2)(1-h/H)$，进度与终局接近程度的交互，用来表达临近终局时
    避免留下未完成进度的倾向。
    \item $\ind\{c_i=1\}$，当前是否处于第二面。
    \item $\ind\{x_i<6\}$，窗户是否仍有效；对所有合法填充动作都等于一，
    因而在动作比较中相当于填充动作的常数偏置。
    \item $\ind\{i=8\}$，是否为最右侧窗户，用来表达普通线性位置项以外的
    边界效应。
    \item $z_i^2$，非线性位置项，允许策略对靠右位置表现出凸或凹的额外偏好。
\end{enumerate}

需要注意，第一版实现中的 $F_i(\xvec)$ 按\emph{行动前}状态计算。如果
$x_i=5$，本次填充会使第 $i$ 扇窗户彻底完成并移出，但 $F_i$ 仍将它计入，
因此会比严格的行动后可达窗户比例高 $1/8$。本节保留这一写法是为了忠实记录
第一版搜索过程；它也说明这些特征是价值代理，而不是精确的Bellman续接价值。

重置另用一个可解释阈值判定。令 $U_{\mathrm{fill}}$ 为当前Mask下最优合法填充
效用，$U_{\mathrm{left}}$ 为左侧未完成窗户的可用概率加权效用，则定义
\[
D_{\mathrm{reset}}
=U_{\mathrm{left}}-U_{\mathrm{fill}}
+\beta_0+\beta_1R_{\mathrm{left}}
+\beta_2\frac{g-1}{7}+\beta_3\frac{h}{H}.
\]
当 $D_{\mathrm{reset}}>0$ 时主动重置；没有合法填充时按照第一章规则被迫重置或
空过。这里 $R_{\mathrm{left}}$ 表示玻璃匠左侧仍然有效的窗户比例。

这种设计的优点是每个动作可以被拆解解释，缺点是填充和重置都依赖人工特征；
线性效用也未必能表达策略阶段切换或多步协同。

\section{参数搜索与评估方法}

参数搜索先进行广域随机搜索，再围绕表现较好的候选使用逐渐减小方差的高斯扰动
做局部细化。为了降低两个候选策略比较时的蒙特卡洛噪声，训练阶段使用共同随机数：
预先固定 $U_{e,t,i}\sim\mathrm{Uniform}(0,1)$，在状态进度为 $q_i$ 时令
\[
m_{e,t,i}=\ind\{U_{e,t,i}<p(q_i)\}.
\]
这样，同一批候选面对同一组底层随机数，参数变化造成的分数差更容易被辨识。

最终报告不复用训练随机数，而使用十万条新模拟轨迹，报告平均分、平均完成面数、
重置次数和窗户级分布。开局动作另外采用单步rollout改进：固定第一步选择某一扇
窗户，之后沿当前策略继续，在配对随机数下比较各第一步的期望结果。需要强调，
这只是对开局的一步策略改进，不等于对所有可达状态递归执行完整策略迭代。

\paragraph{实现校验中的修正。}
早期代码从 $H=36$ 实验沿用了时间特征的旧分母，使本应为 $h/H$ 的特征在
$H=45$ 时没有同步标准化。发现后已统一使用当前的 $H$。因此，修正前曾出现的
约 $36.00$ 分不再作为有效基准；下文所有对比均使用重新评估后的 $35.91$ 分。

\section{第一版结果、讨论反馈与局限诊断}

\paragraph{初始结果与中间结论。}
在 $H=45$、十万条新轨迹上，第一版策略平均得到 $35.9135$ 分（标准误
$0.0103$），平均完成 $9.313$ 面，并使用 $14.075$ 次重置。它通常先让几乎
所有窗户至少完成一面，随后将第二面集中在第6、7、8扇窗户。因此当时形成了
“先完成所有第一面，再主要完成6至8号第二面”的描述。该描述忠实反映第一版
策略的行为，但并不是最优性结论。

\paragraph{讨论反馈与支配关系。}
讨论中提出：既然小号窗户完成第二面时可以再次享受更多右侧加分，为什么策略反而
着重完成大号窗户？若所有右侧窗户都已激活，第1至第8扇窗户每次完成的分数依次为
\[
(8,7,6,5,4,3,2,1).
\]
第二面不会创造新的激活价值。因此，在进度、Mask、移动与重置成本相同的条件下，
小号窗户的第二面严格优于大号窗户。大号窗户只有在更接近完成、位置成本更低、
小号窗户被Mask遮蔽或时域即将结束时，才可能成为局部更优选择。

\paragraph{右侧激活数是否已经进入第一版特征？}
令候选窗户 $i$ 右侧至少完成过一面的窗户数为
\[
R_i(\xvec)=\sum_{j=i+1}^{8}\ind\{x_j\ge3\}.
\]
这个量确实出现在第一版的第5个特征中，因为
\[
\phi_5=\frac{d_iB_i(\xvec)}{8}
=\ind\{q_i=2\}\frac{1+R_i(\xvec)}{8}.
\]
但它\emph{没有作为独立状态特征出现}：只有当本次填充立即完成一面，即
$q_i=2$ 时，策略才能看到 $R_i$；当 $q_i=0$ 或 $1$ 时，$\phi_5=0$，右侧
激活数的未来价值完全消失。这正是第一版策略低估尚需两至三步才能完成的小号窗户
的主要原因，而后续加入的独立 $R_i$ 特征正是对此进行修正。

\paragraph{局限诊断。}
综合结果与上述反馈，第一版策略存在五个主要局限：右侧位置偏置过强；右侧激活
收益只在“本次完成”时显现；没有显式区分“向右建设”和“向左收割”两个阶段；
初始参数只经过暖启动和局部搜索，并非联合全局优化；重置决策比较的是人工效用
代理而不是真实的多步 $Q$ 值。这些问题共同说明，第一版观察到的右侧第二面集中
是策略表示能力造成的行为偏差，而不是由计分规则推出的可靠洞见。

\section{独立右侧激活特征与联合重搜}

为了让右侧结构价值在尚未完成当前面时也保持可见，同时避免预先规定“剩余一步
比剩余三步价值高三倍”，本轮不使用带有剩余填充次数分母的得分密度特征。直接将
\[
\widetilde R_i(\xvec)
=\frac{R_i(\xvec)}{7}
=\frac{1}{7}\sum_{j=i+1}^{8}\ind\{x_j\ge3\}
\]
作为第15个独立特征。除以7只用于数值归一化，不改变窗户之间的顺序。新的动作
效用为
\[
u_i^{R}(s,h)
=\boldsymbol{w}_{1:14}^{\mathsf T}\boldsymbol{\phi}(s,i,h)
+w_R\widetilde R_i(\xvec).
\]
这里，完成进度仍由原特征 $q_i/2$、$\ind\{q_i=2\}$ 及其时域交互项表达；
$\widetilde R_i$ 只表达右侧已经建立的计分结构。把两者拆开后，搜索可以自行学习
结构收益与完成紧迫度之间的权衡，而不是通过 $1/(3-q_i)$ 硬编码该权衡。

本轮使用共同随机数，在六千条训练轨迹上比较四个启发式种子、二百个广域随机
候选和四轮共四百八十个高斯局部扰动。原十四个动作权重、新增的 $w_R$ 和四个
重置参数全部联合调整。最终得到
\[
w_R=14.3267>0,
\]
说明在当前归一化下，独立的右侧激活信息确实被搜索稳定地赋予正价值。训练完成后
固定全部参数，另用十万条完全独立的随机轨迹进行下节验证。

\section{独立右侧激活特征策略的 H=45 结果}

十万条新轨迹的验证结果见表\ref{tab:overall}。

\begin{table}[htbp]
\centering
\caption{第一版策略与独立 $R_i$ 特征策略的独立样本验证}
\label{tab:overall}
\begin{tabular}{lrr}
\toprule
指标 & 第一版策略 & 独立 $R_i$ 策略 \\
\midrule
平均得分 & 35.9135 & 44.2047 \\
得分标准误 & 0.0103 & 0.0169 \\
平均完成面数 & 9.313 & 10.244 \\
平均重置次数 & 14.075 & 11.249 \\
其中：主动重置 & --- & 9.667 \\
其中：被迫重置或空过 & --- & 1.582 \\
\bottomrule
\end{tabular}
\end{table}

新策略的分数第10、50和90百分位数分别为 $37$、$44$ 和 $51$。平均得分比第一版
提高 $8.2912$ 分，即约 $23.1\%$；平均完成面数增加 $0.9306$，平均重置减少
$2.8258$ 次。因此，这一改进同时来自更有价值的完成顺序和更高的行动利用效率，
而不只是把相同数量的完成重新分配位置。

表\ref{tab:window}给出窗户级分布。第6至第8扇窗户几乎只完成第一面，用来建立
右侧加分结构；第二面则集中在第1至第5扇窗户，其中第2至第4扇最常完成两面。
这与讨论中提出的总体方向一致，并推翻了“第二面应集中在6至8号”的中间结论。

\begin{table}[htbp]
\centering
\small
\caption{各窗户完成分布；概率以百分比表示}
\label{tab:window}
\begin{tabular}{crrrrr}
\toprule
窗户 & \multicolumn{2}{c}{第一版策略} & \multicolumn{3}{c}{独立 $R_i$ 策略} \\
\cmidrule(lr){2-3}\cmidrule(lr){4-6}
$i$ & 平均完成面 & 完成两面概率 & 平均完成面 & 完成两面概率 & 至少一面概率 \\
\midrule
1 & 0.897 & 0.0\% & 1.372 & 40.0\% & 97.2\% \\
2 & 0.986 & 0.0\% & 1.499 & 53.0\% & 96.9\% \\
3 & 0.998 & 0.0\% & 1.714 & 71.6\% & 99.8\% \\
4 & 1.000 & 0.0\% & 1.552 & 55.2\% & 100.0\% \\
5 & 1.026 & 2.6\% & 1.103 & 10.3\% & 100.0\% \\
6 & 1.306 & 30.6\% & 1.004 & 0.4\% & 100.0\% \\
7 & 1.784 & 78.4\% & 1.000 & 0.0\% & 100.0\% \\
8 & 1.317 & 31.7\% & 0.999 & 0.0\% & 99.9\% \\
\bottomrule
\end{tabular}
\end{table}

结果并不意味着第二面应严格按窗户编号从小到大完成。第3扇窗户完成两面的概率
高于第1扇，说明右侧激活收益还必须与Mask、已有进度、玻璃匠位置和重置成本共同
权衡。独立 $R_i$ 修正的是“尚未到完成前一步时看不见结构收益”的缺陷，但它本身
并不决定完整的动作排序。

\section{同时保留三类结构特征的更完备搜索}

独立 $R_i$ 将结构收益与进度拆开，但不代表带剩余填充成本的密度信息一定无用。
因此，下一轮同时保留三个可由搜索自由组合的特征。除上一节的
$\widetilde R_i=R_i/7$ 外，定义
\[
\widetilde D_i(s)
=\frac{1+R_i(\xvec)}{8(3-q_i)},
\qquad
L_i(s)
=\frac{A(\xvec)}{8}\left(1-\frac{i-1}{7}\right),
\]
其中
\[
A(\xvec)=\sum_{j=1}^{8}\ind\{x_j\ge3\}
\]
是全局已经至少完成一面的窗户数。$\widetilde R_i$ 表示纯粹的右侧结构收益，
$\widetilde D_i$ 表示结构收益相对于当前面剩余成功填充次数的密度，$L_i$ 则表示
随着全局激活程度提高而增强的左侧收割倾向。完整动作效用为
\[
u_i^{RDL}(s,h)
=\boldsymbol{w}_{1:14}^{\mathsf T}\boldsymbol{\phi}(s,i,h)
+w_R\widetilde R_i(s)+w_D\widetilde D_i(s)+w_L\,L_i(s).
\]
三类特征同时进入模型的意义是：不预先决定“是否应该除以剩余填充次数”，而让
联合搜索从数据中确定三者的边际作用。

本轮从已经验证的 $R_i$ 策略暖启动，使用两个独立随机种子进行完整重启。每个
重启使用五千条共同随机数训练轨迹，并比较四个启发式种子、一个 $R_i$ 暖启动、
四十九组 $(w_D,w_L)$ 结构化网格、百余个广域随机候选以及五轮共五百个局部高斯
扰动。两个重启先在各自独立的一万二千条轨迹上初测，再在同一批十万条新轨迹上
做配对筛选；选定候选后，另用一批从未参与搜索或选择的十万条轨迹作最终报告。

最终参数中三个新增特征的权重为
\[
(w_R,w_D,w_L)=(15.3966,-7.0428,4.2708).
\]
$w_R>0$ 和 $w_L>0$ 表明右侧结构收益与阶段化左侧收割均被保留。$w_D<0$ 并不
等于删除 $D_i$：它表示在已经控制 $R_i$、$L_i$、当前进度、立即完成指标和时间
紧迫度后，再额外正向奖励“单位剩余填充成本得分”会重复偏向临门一脚，搜索反而
需要负系数校正这一倾向。由于这些特征高度相关，系数符号应被解释为联合模型中的
边际作用，而不能单独作因果解释。

\section{三类结构特征联合策略的 H=45 结果}

表\ref{tab:rdl-overall}比较第一版、独立 $R_i$ 策略与最终 $RDL$ 策略。前两列
沿用各自原始的十万样本报告，最终列来自完全未参与候选筛选的十万条轨迹。

\begin{table}[htbp]
\centering
\caption{三轮策略的独立样本验证结果}
\label{tab:rdl-overall}
\begin{tabular}{lrrr}
\toprule
指标 & 第一版策略 & 独立 $R_i$ 策略 & $RDL$ 联合策略 \\
\midrule
平均得分 & 35.9135 & 44.2047 & 45.0538 \\
得分标准误 & 0.0103 & 0.0169 & 0.0170 \\
平均完成面数 & 9.313 & 10.244 & 10.257 \\
平均重置次数 & 14.075 & 11.249 & 11.061 \\
其中：主动重置 & --- & 9.667 & 9.469 \\
其中：被迫重置或空过 & --- & 1.582 & 1.591 \\
\bottomrule
\end{tabular}
\end{table}

最终策略的分数第10、50和90百分位数分别为 $38$、$45$ 和 $52$。在最终测试所用
的同一批随机轨迹上，$R_i$ 策略平均得到 $44.2219$ 分，$RDL$ 策略的配对提升为
\[
0.8318\pm0.0202,
\]
这里误差项是配对差值的标准误。提升约为四十一个标准误，说明它不是普通模拟
噪声。与此同时，完成面数只增加 $0.0101\pm0.0031$，所以这约 $0.83$ 分的收益
几乎全部来自更好的完成位置与时序，而不是做出更多完成。

相对最初的十四特征策略，最终平均得分提高 $9.1403$ 分，约为 $25.5\%$；平均
完成面数增加约 $0.944$，平均重置减少约 $3.014$ 次。表\ref{tab:rdl-window}
进一步给出最终策略的窗户级行为。

\begin{table}[htbp]
\centering
\caption{$RDL$ 联合策略的窗户级完成分布}
\label{tab:rdl-window}
\begin{tabular}{crrr}
\toprule
窗户 $i$ & 平均完成面 & 至少完成一面概率 & 完成两面概率 \\
\midrule
1 & 1.274 & 92.7\% & 34.7\% \\
2 & 1.542 & 97.6\% & 56.6\% \\
3 & 1.764 & 99.7\% & 76.6\% \\
4 & 1.581 & 100.0\% & 58.1\% \\
5 & 1.094 & 100.0\% & 9.4\% \\
6 & 1.004 & 100.0\% & 0.4\% \\
7 & 1.000 & 100.0\% & 0.0\% \\
8 & 0.998 & 99.8\% & 0.0\% \\
\bottomrule
\end{tabular}
\end{table}

最终策略仍把第6至第8扇窗户主要作为一次性结构激活器，并把第二面集中在第1至
第5扇窗户。第3、4、2扇的第二面概率最高，再次说明“越靠左越好”只是条件相同时
的支配关系；实际排序还要支付玻璃匠移动、重置、Mask与已有进度的机会成本。

\section{高权重特征的二阶扩展与筛选}

前述 $RDL$ 策略仍然是线性效用：某一特征的边际作用不随其他特征改变。为了检验
“进度、第二面与右侧结构是否应当相互调节”，本轮从最终一阶权重绝对值较大的
六项出发，定义
\[
P_i=\frac{q_i}{2},\qquad
C_i=\ind\{q_i=2\},\qquad
S_i=\ind\{c_i=1\},
\]
并与 $\widetilde R_i,\widetilde D_i,L_i$ 合并成
\[
\boldsymbol\psi_i=(P_i,C_i,S_i,\widetilde R_i,\widetilde D_i,L_i)^{\mathsf T}.
\]
所谓二阶项是特征乘积 $\psi_{ia}\psi_{ib}$，每个乘积拥有独立待搜索系数；它并
不是把两个一阶权重直接相乘。所有父项仍保留在模型中，因此这是一个层级式二阶
模型。

六项可产生十五个两两交互和六个平方项，但其中三项没有增加新的表示能力：
\[
P_iC_i=C_i,\qquad C_i^2=C_i,\qquad S_i^2=S_i.
\]
第一式成立是因为 $C_i=1$ 时必有 $P_i=1$；后两式来自指标变量的二元性。因此，
真正需要验证的是十四个非重复交互，以及
\[
P_i^2,\quad \widetilde R_i^2,\quad \widetilde D_i^2,\quad L_i^2
\]
四个非重复平方项。

筛选阶段固定已经验证的 $RDL$ 参数，在同一批一万六千条共同随机数轨迹上，逐项
尝试系数网格
\[
\{-32,-20,-12,-6,0,6,12,20,32\}.
\]
表\ref{tab:interaction-screen}给出每项单独加入时的最优网格系数与相对 $RDL$
的配对平均分增量。这里的系数只用于筛选，不是最终联合模型的估计值。

\begin{table}[htbp]
\centering
\scriptsize
\caption{二阶候选的单项共同随机数筛选；增量误差为配对标准误}
\label{tab:interaction-screen}
\begin{tabular}{lrrr@{\qquad}lrrr}
\toprule
交互项 & 系数 & 增量 & 标准误 & 交互项 & 系数 & 增量 & 标准误 \\
\midrule
$P_iS_i$ & 6 & 0.1017 & 0.0336 & $P_i\widetilde R_i$ & 6 & 0.0130 & 0.0396 \\
$P_i\widetilde D_i$ & 0 & 0 & 0 & $P_iL_i$ & 6 & 0.0295 & 0.0392 \\
$C_iS_i$ & 6 & 0.0604 & 0.0213 & $C_i\widetilde R_i$ & 0 & 0 & 0 \\
$C_i\widetilde D_i$ & 0 & 0 & 0 & $C_iL_i$ & 0 & 0 & 0 \\
$S_i\widetilde R_i$ & 0 & 0 & 0 & $S_i\widetilde D_i$ & 6 & 0.1145 & 0.0362 \\
$S_iL_i$ & 6 & 0.0049 & 0.0434 & $\widetilde R_i\widetilde D_i$ & 6 & 0.1460 & 0.0305 \\
$\widetilde R_iL_i$ & 0 & 0 & 0 & $\widetilde D_iL_i$ & 0 & 0 & 0 \\
\midrule
平方项 & 系数 & 增量 & 标准误 & 平方项 & 系数 & 增量 & 标准误 \\
\midrule
$P_i^2$ & 0 & 0 & 0 & $\widetilde R_i^2$ & 0 & 0 & 0 \\
$\widetilde D_i^2$ & 6 & 0.0402 & 0.0321 & $L_i^2$ & 0 & 0 & 0 \\
\bottomrule
\end{tabular}
\end{table}

信号最强的四个交互依次为
\[
\widetilde R_i\widetilde D_i,\quad
S_i\widetilde D_i,\quad
P_iS_i,\quad C_iS_i.
\]
前三项的筛选增量约为三个至五个配对标准误，$C_iS_i$ 也接近三个标准误。
$\widetilde D_i^2$ 的单项信号较弱，只有约 $1.25$ 个标准误；但为检验密度特征
是否存在曲率，本轮仍将它保留到联合重搜。其余候选固定为零。这种先筛选再联合
估计的做法控制了参数数量，但也意味着本轮没有穷举原十七个一阶特征之间的全部
交互。

\section{二阶联合重搜与独立验证}

联合阶段同时调整十七个一阶动作权重、上述四个交互权重、$\widetilde D_i^2$
权重和四个重置参数。每个独立重启使用四千五百条共同随机数训练轨迹，比较七个
暖启动、一百二十个广域候选和五层共三百五十个局部高斯扰动，总计四百七十七次
策略评估。另跑一个不含平方项的纯交互重启作为消融对照。

两个完整重启和纯交互候选随后在同一批十万条新轨迹上筛选，结果见
表\ref{tab:second-order-selection}。纯交互候选虽然平均完成面数提高到 $10.570$，
得分却几乎没有变化；这再次说明完成数量不是充分目标，错误位置上的第二面可能
消耗行动却只获得较低分。允许 $\widetilde D_i^2$ 并联合重估全部父项后，两个独立
重启都产生稳定正增量，第二个重启略优并被选中。不过，纯交互与完整模型来自不同
搜索路径，二者差异不能全部归因于平方项；后文另作固定参数消融。

\begin{table}[htbp]
\centering
\caption{十万条共同选择轨迹上的二阶候选比较}
\label{tab:second-order-selection}
\begin{tabular}{lrrrr}
\toprule
策略 & 平均分 & 对 $RDL$ 增量 & 增量标准误 & 完成面数 \\
\midrule
$RDL$ 基准 & 45.0596 & --- & --- & 10.256 \\
纯交互重启 & 45.0627 & 0.0031 & 0.0196 & 10.570 \\
完整二阶重启1 & 45.2366 & 0.1771 & 0.0160 & 10.270 \\
完整二阶重启2 & 45.2695 & 0.2099 & 0.0160 & 10.278 \\
\bottomrule
\end{tabular}
\end{table}

被选策略的五个二阶系数为
\[
\begin{aligned}
(&w_{PS},w_{CS},w_{S\widetilde D},
w_{\widetilde R\widetilde D},w_{\widetilde D^2})\\
&=(0.8604,\ 0.5630,\ -0.2047,\ 5.1079,\ -0.1879).
\end{aligned}
\]
筛选阶段和联合阶段的系数符号不必一致：筛选时其余参数固定，联合阶段则允许
高度相关的一阶父项与二阶项相互替代。可靠的结论是完整策略在新轨迹上的行为与
收益，而不是把单个系数解释成独立因果效应。数值上，最大的新增项是
$\widetilde R_i\widetilde D_i$：当右侧激活结构与单位剩余填充价值同时较高时，
策略给予额外奖励，这正是线性 $RDL$ 无法表达的协同。

候选选定后，再使用一批从未参与训练或候选选择的十万条轨迹作最终报告。结果见
表\ref{tab:second-order-final}。二阶策略平均得到 $45.2569$ 分；在同一批轨迹上，
$RDL$ 得到 $45.0644$ 分，配对提升为
\[
0.1925\pm0.0160.
\]
提升约为十二个配对标准误，因此不是普通模拟噪声。完成面数只增加
$0.0195\pm0.0024$，而重置次数增加 $0.2237\pm0.0040$；新增收益主要来自愿意
支付少量额外重置以完成更高价值的左侧窗口，而不是单纯提高行动利用率。

\begin{table}[htbp]
\centering
\caption{最终未触碰十万条轨迹上的 $RDL$ 与二阶策略}
\label{tab:second-order-final}
\begin{tabular}{lrr}
\toprule
指标 & $RDL$ 策略 & 二阶策略 \\
\midrule
平均得分 & 45.0644 & 45.2569 \\
得分标准误 & 0.0170 & 0.0168 \\
平均完成面数 & 10.2585 & 10.2780 \\
平均重置次数 & 11.0480 & 11.2717 \\
得分第10/50/90百分位 & 38/45/52 & 38/45/52 \\
\bottomrule
\end{tabular}
\end{table}

为单独判断平方项的边际贡献，在最终参数中仅将
$w_{\widetilde D^2}$ 从 $-0.1879$ 设为零，其余一阶、交互和重置参数完全不变。
同一批十万条轨迹上，得分从 $45.2569$ 降至 $45.2509$，即平方项贡献
\[
0.0061\pm0.0022.
\]
因此 $\widetilde D_i^2$ 确有小幅正增量，但只解释总提升的一小部分；主要收益来自
交互结构与全部父项的联合重估。若要进一步确认平方项的稳定性，还应分别以“固定
不含平方项”和“允许平方项”为条件进行更多成对搜索重启，而不能只比较一次最优
候选。

窗户级结果见表\ref{tab:second-order-window}。相较 $RDL$，二阶策略进一步把第二面
从第4扇重新分配到第1、2扇；第6至第8扇仍主要作为一次性结构激活器。因此，约
$0.19$ 分的提升具有清楚的行为解释：不是改变“先激活右侧、再收割左侧”的总体
框架，而是在左侧收割阶段更精细地利用进度、第二面与密度的联合条件。

\begin{table}[htbp]
\centering
\caption{最终二阶策略的窗户级完成分布}
\label{tab:second-order-window}
\begin{tabular}{crrr}
\toprule
窗户 $i$ & 平均完成面 & 至少完成一面概率 & 完成两面概率 \\
\midrule
1 & 1.361 & 94.4\% & 41.7\% \\
2 & 1.643 & 98.1\% & 66.2\% \\
3 & 1.760 & 99.8\% & 76.2\% \\
4 & 1.435 & 100.0\% & 43.5\% \\
5 & 1.079 & 100.0\% & 7.9\% \\
6 & 1.006 & 100.0\% & 0.6\% \\
7 & 1.000 & 100.0\% & 0.0\% \\
8 & 0.994 & 99.4\% & 0.0\% \\
\bottomrule
\end{tabular}
\end{table}

\section{补齐重置特征空间后的配对Rollout重搜}

前一版学习型重置策略只使用二十个手工汇总特征。虽然它已经优于四参数线性阈值，
但与填充选点策略的三十六维动作表示并不对称：重置模型只能看到少量密度、位置和
效用差，无法完整辨认当前准备放弃的填充，也无法描述重置后重新开放的整组左侧
候选。这个限制会把不同盘面压缩成相同的重置状态。

令 $\phi_i(s)\in\mathbb{R}^{36}$ 为上一节最终填充策略对窗户 $i$ 使用的完整动作
特征，其中包含十七个基础量、十五个两两交互和四个平方项。令 $a$ 为当前Mask下
得分最高的合法填充，$\mathcal L(s)=\{j<g:x_j<6\}$ 为重置后重新开放的左侧窗户，
并定义
\[
j^*=\arg\max_{j\in\mathcal L(s)}p_j U_j(s),
\]
其中 $p_j\in\{0.8,0.6,0.4\}$ 是该窗下一次填充在新Mask中可用的概率。扩展后的
重置表示为
\[
z(s)=\Big[z^{(20)}(s),\ \phi_a,\ p_{j^*}\phi_{j^*},\
p_{j^*}\phi_{j^*}-\phi_a,\
\operatorname{mean}_{j\in\mathcal L}p_j\phi_j,\
\operatorname{max}_{j\in\mathcal L}p_j\phi_j\Big]\in\mathbb{R}^{200},
\]
其中最后一个最大值逐坐标计算。维数正好为 $20+5\times36=200$。这一表示同时
覆盖“当前填什么”“重置最可能换来什么”以及“重置打开的候选池是否有冗余”，并
继承填充策略已经验证过的一阶、交互和曲率结构。

标签通过共同随机数配对rollout生成。在旧二阶策略访问的一万局中共观察到
$266{,}728$ 个可选重置状态，从中不放回抽取 $1{,}600$ 个；对每个状态分别强制
“立即重置”和“执行当前最佳填充”，再共享 $128$ 组后续Mask直到 $H=45$。监督
目标为两种动作的后续总分差。训练集占 $80\%$，其余 $20\%$ 只用于选择回归器
结构；随后另用三万局搜索决策阈值。为公平检验特征扩展，同一标签、同一切分上
另训练只读取前二十列的HGB消融模型。

最佳二百维模型使用 $15$ 个叶节点、最小叶样本数 $20$ 和 $L_2=2$。验证结果见
表\ref{tab:reset-feature-ablation}。扩展模型把RMSE从 $0.6539$ 降到 $0.6249$，
并把重置收益正负号判断准确率从 $91.56\%$ 提高到 $92.19\%$。三万局端到端选择
集中，二十维消融模型的最佳点为阈值 $0.05$、平均分 $45.5838$；二百维模型在
阈值 $0.10$ 达到 $45.6209$。因此增益并非单纯来自增加树模型容量，而确实依赖
新增状态信息。

\begin{table}[htbp]
\centering
\caption{相同配对标签上的重置特征空间消融}
\label{tab:reset-feature-ablation}
\begin{tabular}{lrrrr}
\toprule
模型 & 特征数 & 验证RMSE & 符号准确率 & 三万局最佳分 \\
\midrule
旧HGB消融 & 20 & 0.6539 & 91.56\% & 45.5838 \\
扩展HGB & 200 & 0.6249 & 92.19\% & 45.6209 \\
\bottomrule
\end{tabular}
\end{table}

最后在完全未参与训练和阈值选择的十万局上评估。表
\ref{tab:expanded-reset-final}中的增量均与旧四参数重置规则使用同一组Mask配对
计算。阈值 $0.10$ 的期望分最高；阈值 $0.20$ 只低 $0.0081$ 分，远小于单个均值
标准误，却再少约 $0.172$ 次重置并多完成 $0.034$ 面，因此作为默认平衡推荐。
$0.35$ 适合更重视减少重置的目标；到 $0.50$ 时已经越过有效前沿并开始损失分数。

\begin{table}[htbp]
\centering
\caption{二百维重置策略在最终十万条未触碰轨迹上的结果}
\label{tab:expanded-reset-final}
\begin{tabular}{lrrrrr}
\toprule
策略 & 平均分 & 得分标准误 & 完成面数 & 总重置 & 配对增分 \\
\midrule
旧四参数重置 & 45.2889 & 0.0168 & 10.2813 & 11.2767 & --- \\
阈值0.10 & 45.5747 & 0.0167 & 10.4562 & 10.6447 & 0.2858 \\
阈值0.20（推荐） & 45.5666 & 0.0168 & 10.4906 & 10.4726 & 0.2777 \\
阈值0.35 & 45.4993 & 0.0168 & 10.5340 & 10.2962 & 0.2104 \\
阈值0.50 & 45.2379 & 0.0169 & 10.5629 & 10.0623 & -0.0510 \\
\bottomrule
\end{tabular}
\end{table}

以推荐阈值 $0.20$ 为例，可选重置从 $9.6952$ 降为 $8.8667$，被迫重置从
$1.5815$ 小幅增至 $1.6059$；总重置净减少 $0.8041$ 次。配对分数增量为
$0.2777\pm0.0176$，完成面增量为 $0.2093\pm0.0026$。这说明旧策略的主要问题
不是普遍“太保守”或“太激进”，而是无法区分哪些左侧候选池值得花一整个行动
重新开放。扩展模型学到的是状态依赖的选择性重置。

本章模拟器、二阶搜索与重置rollout脚本已发布在
\href{https://github.com/CYZ-gmail/azul-stained-glass-strategy}{公开GitHub仓库}。
上述实验对应的源代码固定快照为
\href{https://github.com/CYZ-gmail/azul-stained-glass-strategy/commit/550f1d63434b215ae0770dcfb57493ef20f6c063}%
{commit \texttt{550f1d6}}。Overleaf编译不依赖远程分支内容，因此文中使用不可变
commit链接而不是动态读取GitHub；这样即使后续继续修改主分支，当前数字仍能追溯
到准确脚本版本。

\section{当前方法的边界与下一步改进}

截至本章，结论应被解释为“在当前可解释策略族中发现的高质量策略”，而不是全局
最优策略。主要局限与对应改进方向如下。

\begin{enumerate}[leftmargin=*]
    \item \textbf{策略族限制：}当前虽加入五个筛选后的二阶项，但填充动作打分仍是
    固定基函数的线性组合；重置已改为二百维非线性模型，但只完成一次策略改进。
    下一步可做近似策略迭代：采样高频访问状态，对所有合法动作做配对rollout，
    用改进后的动作标签拟合决策树或非线性价值模型，再重复迭代。
    \item \textbf{参数搜索限制：}十七个一阶权重、五个二阶权重与四个重置参数
    已经联合重搜，并使用两个独立搜索种子和最终未触碰随机数验证；但随机搜索加
    局部扰动仍不能保证找到该策略族的全局最优参数。候选筛选也只覆盖六个高权重
    一阶特征，而不是全部特征的二阶展开。下一轮可增加搜索重启，或改用近似策略
    迭代。
    \item \textbf{重置标签的策略不一致：}配对rollout已经直接比较“本回合重置”
    与“当前填充”的多步 $Q$ 值，但两条分支后续仍由旧策略接管，因此只是一次
    近似策略改进。下一轮应让新重置策略重新收集访问状态并迭代拟合，直到分布与
    决策基本稳定；还应使用单独测试集直接比较二十维和二百维模型。
    \item \textbf{开局结论尚需重算：}此前的开局rollout基于旧策略；收割逻辑修正
    后必须重新计算八个强制首步的动作价值，才能回答最终的开局优先级。
    \item \textbf{随机性简化：}独立Mask忽略了共享花色供给造成的相关性。后续可
    加入公共“场面质量”变量或显式颜色组成，对比独立与相关Mask。
    \item \textbf{规则抽象：}当前没有显式建模花色、对手、碎玻璃轨、印刷窗户
    分值以及残留进度终局价值。这些变量都可能改变重置阈值与收割顺序。
\end{enumerate}

建议的下一阶段不是立即增加全部真实规则，而是先在当前模型内完成一次近似策略
迭代，并重算开局动作价值；随后再依次改变 $H\in\{36,40,45,50\}$、Mask概率、
Mask相关性和终局价值。这样可以区分“由算法缺陷造成的行为”与“由规则参数改变
造成的策略变化”。若希望把平均完成面数至少九面作为硬性要求，也可以把搜索目标
写成
\[
\max_{\pi}\;\mathbb{E}_{\pi}[\text{得分}]
\quad\text{s.t.}\quad
\mathbb{E}_{\pi}[\text{完成面数}]\ge 9.
\]
当前推荐策略的平均完成面数为 $10.491$，满足该约束。

\end{document}
